\chapter{Chiral Axial Anomalies}
\label{ch:volume-ii-20}

\section{Classical Chiral Currents}

Let \(\psi\) be a Dirac fermion in four dimensions, coupled to a background or
dynamical gauge field through
\[
  D_\mu\psi=\partial_\mu\psi-\ii A_\mu\psi .
\]
For a mass \(m\), the Dirac Lagrangian is
\[
  \mathcal L
  =
  -\bar\psi\gamma^\mu D_\mu\psi
  -m\bar\psi\psi .
\]
When \(m=0\), the classical equations of motion imply conservation of the
vector and axial currents
\[
  j^\mu=\bar\psi\gamma^\mu\psi,
  \qquad
  j_5^\mu=\bar\psi\gamma^\mu\gamma_5\psi,
  \qquad
  \partial_\mu j^\mu=0,
  \qquad
  \partial_\mu j_5^\mu=0 .
\]
For \(m\neq0\), the axial divergence becomes
\[
  \partial_\mu j_5^\mu=2m\,\bar\psi\gamma_5\psi
\]
with the \(\gamma_5\) convention fixed in the preceding QCD chapter.

The corresponding infinitesimal transformations are
\[
  \delta_V\psi=\ii\alpha\psi,\qquad
  \delta_A\psi=\ii\beta\gamma_5\psi,
\]
where \(\alpha\) and \(\beta\) are constant real parameters.  Classically, for
massless fermions, both transformations are symmetries of the local
Lagrangian.  Quantum mechanically, the vector transformation is preserved
when it is gauged consistently, while the axial current acquires a local
divergence determined by the gauge background.

\section{The Abelian Axial Anomaly}

For a Dirac fermion of electric charge \(q\) coupled to an electromagnetic
field with coupling \(e\), the renormalized axial current may be chosen so
that
\begin{equation}
  \partial_\mu j_5^\mu
  =
  2m\,\bar\psi\gamma_5\psi
  +
  \frac{e^2 q^2}{16\pi^2}
  \epsilon^{\mu\nu\rho\sigma}
  F_{\mu\nu}F_{\rho\sigma}.
  \label{eq:abelian-axial-anomaly}
\end{equation}
The coefficient is fixed by the short-distance singularity of the triangle
diagram with one axial-current insertion and two vector-current insertions.
The vector Ward identity is imposed because the vector current is the current
coupled to the gauge field.  The remaining local term appears in the axial
Ward identity.

Equation \eqref{eq:abelian-axial-anomaly} is a statement about a
renormalized composite operator identity.  Each symbol in it has a prescribed
meaning: \(F_{\mu\nu}\) is the electromagnetic field strength,
\(\epsilon^{0123}=+1\) in the chosen orientation, \(m\bar\psi\gamma_5\psi\)
is the pseudoscalar mass operator, and \(j_5^\mu\) is the renormalized current
whose vector-current Ward identities have been preserved.

\section{Measure Derivation}

The same coefficient is obtained from the path-integral measure.  Work in
Euclidean signature with Dirac operator
\[
  \mathcal D=\gamma^\mu(\partial_\mu-\ii A_\mu).
\]
Expand the fermion fields in an orthonormal basis of eigenfunctions of
\(\mathcal D\).  Under a local chiral rotation
\[
  \psi(x)\mapsto \ee^{\ii\beta(x)\gamma_5}\psi(x),
  \qquad
  \bar\psi(x)\mapsto \bar\psi(x)\ee^{\ii\beta(x)\gamma_5},
\]
the formal measure changes by a Jacobian.  The regulated logarithm of the
Jacobian is
\begin{equation}
  \log J
  =
  2\ii\lim_{M\to\infty}
  \operatorname{Tr}
  \left[
  \beta(x)\gamma_5
  \exp\left(-\mathcal D^2/M^2\right)
  \right],
  \label{eq:fujikawa-regulated-trace}
\end{equation}
where \(\operatorname{Tr}\) is the trace over spacetime, spin, and internal
indices.  The heat-kernel expansion gives the local limit
\begin{equation}
  \lim_{M\to\infty}
  \operatorname{tr}_{\mathrm{spin}}
  \left[
  \gamma_5
  \exp\left(-\mathcal D^2/M^2\right)
  \right](x,x)
  =
  \frac{1}{32\pi^2}
  \epsilon^{\mu\nu\rho\sigma}
  \operatorname{tr}_{R}(F_{\mu\nu}F_{\rho\sigma}),
  \label{eq:heat-kernel-anomaly-density}
\end{equation}
with \(\operatorname{tr}_{R}\) the trace in the fermion representation.  The
anomaly is therefore the regulated trace of \(\gamma_5\) over an
infinite-dimensional fermion space.  Finite-dimensional cyclic trace
manipulations do not remove \eqref{eq:heat-kernel-anomaly-density}, because
the regulator is part of the definition of the composite current.

For compact Euclidean four-manifolds and suitable boundary conditions, the
integrated version is the index theorem
\begin{equation}
  \operatorname{index}\mathcal D
  =
  n_+-n_-
  =
  \frac1{32\pi^2}
  \int
  \epsilon^{\mu\nu\rho\sigma}
  \operatorname{tr}_{R}(F_{\mu\nu}F_{\rho\sigma})\,\dd^4x ,
  \label{eq:dirac-index-anomaly}
\end{equation}
where \(n_\pm\) are the numbers of zero modes of chirality \(\pm1\).

\section{Nonabelian Flavor and Gauge Anomalies}

For several Weyl fermions, anomalies are encoded by traces over their
representations.  Let \(T^A\) be a generator of a global symmetry current and
let \(t^a,t^b\) be generators of gauge or background gauge fields.  The
coefficient of the mixed anomaly is
\begin{equation}
  \mathcal A^{Aab}
  =
  \operatorname{tr}_{\mathrm{left}}
  \bigl(T^A\{t^a,t^b\}\bigr)
  -
  \operatorname{tr}_{\mathrm{right}}
  \bigl(T^A\{t^a,t^b\}\bigr).
  \label{eq:flavor-anomaly-coefficient}
\end{equation}
With background field strengths \(F^a_{\mu\nu}\), the corresponding current
identity has the form
\begin{equation}
  \partial_\mu J_A^\mu
  =
  \frac1{32\pi^2}
  \mathcal A^{Aab}
  \epsilon^{\mu\nu\rho\sigma}
  F^a_{\mu\nu}F^b_{\rho\sigma}
  +
  \text{explicit breaking terms}.
  \label{eq:nonabelian-flavor-anomaly}
\end{equation}

If the anomalous current is the current of a gauge symmetry, the theory is
inconsistent unless the anomaly cancels.  For left-handed Weyl fermions in
representations \(R_i\) of a gauge group, the local cubic gauge-anomaly
condition is
\begin{equation}
  \sum_i
  \operatorname{tr}_{R_i}
  \bigl(t^a\{t^b,t^c\}\bigr)
  =
  0
  \qquad
  \text{for all }a,b,c.
  \label{eq:gauge-anomaly-cancellation}
\end{equation}
This condition is the ghost-number-one obstruction described by BRST
cohomology in Chapter~\ref{ch:volume-ii-18}.  When it holds, the
Slavnov--Taylor identities can be maintained by local counterterms.  When it
fails, no choice of local counterterms gives a unitary gauge theory with the
declared gauge symmetry.

\section{The Strong CP Parameter}

In QCD with \(N_f\) quark flavors, the Euclidean gauge action may include the
topological term
\[
  \ii\theta\,
  \frac1{32\pi^2}
  \int
  \epsilon^{\mu\nu\rho\sigma}
  \operatorname{tr}(F_{\mu\nu}F_{\rho\sigma})\,\dd^4x .
\]
The complex quark mass matrix \(M\) was introduced in
\eqref{eq:qcd-chiral-mass-lagrangian}.  A flavor-singlet axial rotation
changes the phase of \(\det M\) and, by the anomaly, shifts \(\theta\).  With
the convention used in this chapter, the invariant CP-odd parameter is
\begin{equation}
  \bar\theta
  =
  \theta+\arg\det M .
  \label{eq:qcd-physical-theta}
\end{equation}
If the sign convention for \(\gamma_5\) or for the topological density is
reversed, both transformation laws change together and the invariant
combination is correspondingly rewritten.  The invariant statement is that
only the anomaly-invariant combination of the topological angle and the quark
mass phase is physically meaningful.

\section{Anomalous Pion Couplings}

The axial anomaly has observable low-energy consequences.  In two-flavor QCD,
the neutral pion couples to the electromagnetic topological density through
the anomalous divergence of the axial isospin current:
\[
  \partial_\mu J_{5,3}^\mu
  =
  \frac{e^2}{16\pi^2}
  \operatorname{tr}\bigl(T_3 Q^2\bigr)
  \epsilon^{\mu\nu\rho\sigma}
  F_{\mu\nu}F_{\rho\sigma}
  +
  \text{mass terms}.
\]
Here \(Q\) is the quark charge matrix and \(T_3\) is the third generator of
flavor \(SU(2)\).  The same anomaly is reproduced in the pion effective theory
by the Wess--Zumino--Witten term introduced in the next chapter.
